\scanpage{102}
% archive foliation: 101
rev. étale galoisien \inserted{de $\mathbb P^1_T$} engendré par $f$ est de type $\alpha$. Il est évidemment stable par $\gamma_{\mathcal C_0}$, et correspond à un sous-\deleted{foncteur}\inserted{schéma} ouvert et fermé $M_{\mathcal C_0}^{\alpha}$ de $M_{\mathcal C_0}$, stable par $\gamma_{\mathcal C_0}$.

[Il est vide sauf si le cardinal $N$ de $\alpha$ est un multiple commun des $p(s),q(s),r(s)$ et de $\delta$; \deleted{$\unsure{\exists}$} et $M$ est non vide ssi \deleted{$\exists$} $G$ \inserted{de type $\alpha$} est engendrable par des gén.\ \deleted{$\unsure{k,l}$} $\rho_s,\inserted{\sigma,}\rho_f$ satisfaisant $\rho_s\rho_f\sigma=1$, $\sigma^2=1$, et si \deleted{$\unsure{\forall}$} $\exists\ H$ ss-groupe de $G$, avec $\bigcap\text{conjugués de }H=\{e\}$, et $\mathcal C_0$ déduit de la carte combinatoire orientée $E=G/H$ avec les \unsure{op.}\ $\rho_s,\rho_f,\sigma$ \unsure{opérant dessus}]. Mais il est important de réaliser que \unsure{lors même} qu'on aurait $N=\delta$ i.e.\ $H=\{e\}$, la classe d'isom. de $G$ \uline{muni de} $\rho_s,\rho_f,\sigma$ (i.e. le type de carte combinatoire orientée correspondante) \uline{n'est pas} déterminée par la seule connaissance de $X\xrightarrow{f}\mathbb P^1_{\mathbb C}$\ldots]
